\subsection{The twenty-three lambencies, computed as covers}

\begin{result}{D.14}{Proved}{All passports, cusp fields, Atkin--Lehner actions, and signatures}
For each lambency in Cheng--Duncan--Harvey's Table~2, the following tables give the complete permutation cover $X_0(N)\to X(1)$, its cusps, the action of the subgroup $S$ of exact divisors, and the signature of $\Gamma_0(N)+S$. The signature convention is $(0;e_2,e_3;c)$: genus zero, $e_i$ elliptic points of order $i$, and $c$ cusps. The genus $g_0$ in the first table is that of $X_0(N)$ before taking the $S$-quotient. A cycle expression $a^b$ means $b$ cycles of length $a$.

These are signatures of the effective modular orbifold. The homomorphism $\mathrm{SL}_2(\mathbb Z)\to\mathrm{PSL}_2(\mathbb Z)$ has kernel $\{\pm I\}$, which acts trivially on every dart. Pulling the effective quotient back through this central extension restores the moduli stack of elliptic curves: its generic inertia is $\mu_2$ and its elliptic inertia orders are $4$ and $6$. Thus the displayed finite cover and its unrigidified modular stack are related by that specified central gerbe.
\end{result}

{\small\begin{longtable}{r|r|r|l|l|l|l}
$\ell$ & $\mu$ & $g_0$ & $\sigma_2$ & $\sigma_3$ & cusp widths & $(e_2,e_3;c)$\\\hline
$2$ & 3 & 0 & $1 2$ & $3$ & $1 2$ & $(1,0;2)$\\
$3$ & 4 & 0 & $2^{2}$ & $1 3$ & $1 3$ & $(0,1;2)$\\
$4$ & 6 & 0 & $2^{3}$ & $3^{2}$ & $1^{2} 4$ & $(0,0;3)$\\
$5$ & 6 & 0 & $1^{2} 2^{2}$ & $3^{2}$ & $1 5$ & $(2,0;2)$\\
$6$ & 12 & 0 & $2^{6}$ & $3^{4}$ & $1 2 3 6$ & $(0,0;4)$\\
$7$ & 8 & 0 & $2^{4}$ & $1^{2} 3^{2}$ & $1 7$ & $(0,2;2)$\\
$8$ & 12 & 0 & $2^{6}$ & $3^{4}$ & $1^{2} 2 8$ & $(0,0;4)$\\
$9$ & 12 & 0 & $2^{6}$ & $3^{4}$ & $1^{3} 9$ & $(0,0;4)$\\
$10$ & 18 & 0 & $1^{2} 2^{8}$ & $3^{6}$ & $1 2 5 10$ & $(2,0;4)$\\
$12$ & 24 & 0 & $2^{12}$ & $3^{8}$ & $1^{2} 3^{2} 4 12$ & $(0,0;6)$\\
$13$ & 14 & 0 & $1^{2} 2^{6}$ & $1^{2} 3^{4}$ & $1 13$ & $(2,2;2)$\\
$16$ & 24 & 0 & $2^{12}$ & $3^{8}$ & $1^{4} 4 16$ & $(0,0;6)$\\
$18$ & 36 & 0 & $2^{18}$ & $3^{12}$ & $1^{3} 2^{3} 9 18$ & $(0,0;8)$\\
$25$ & 30 & 0 & $1^{2} 2^{14}$ & $3^{10}$ & $1^{5} 25$ & $(2,0;6)$\\
$6+3$ & 12 & 0 & $2^{6}$ & $3^{4}$ & $1 2 3 6$ & $(2,0;2)$\\
$10+5$ & 18 & 0 & $1^{2} 2^{8}$ & $3^{6}$ & $1 2 5 10$ & $(3,0;2)$\\
$14+7$ & 24 & 1 & $2^{12}$ & $3^{8}$ & $1 2 7 14$ & $(4,0;2)$\\
$18+9$ & 36 & 0 & $2^{18}$ & $3^{12}$ & $1^{3} 2^{3} 9 18$ & $(2,0;4)$\\
$22+11$ & 36 & 2 & $2^{18}$ & $3^{12}$ & $1 2 11 22$ & $(6,0;2)$\\
$30+15$ & 72 & 3 & $2^{36}$ & $3^{24}$ & $1 2 3 5 6 10 15 30$ & $(8,0;4)$\\
$46+23$ & 72 & 5 & $2^{36}$ & $3^{24}$ & $1 2 23 46$ & $(12,0;2)$\\
$12+4$ & 24 & 0 & $2^{12}$ & $3^{8}$ & $1^{2} 3^{2} 4 12$ & $(0,0;4)$\\
$30+6,10,15$ & 72 & 3 & $2^{36}$ & $3^{24}$ & $1 2 3 5 6 10 15 30$ & $(6,0;2)$\\
\end{longtable}
}

\begin{proof}
Use primitive row vectors $(x,y)$ modulo $N$, modulo multiplication by a unit. These are the cosets $\Gamma_0(N)\backslash\mathrm{SL}_2(\mathbb Z)$, by sending a matrix to its bottom row. Its stabilizer is exactly $\Gamma_0(N)$; lifting a primitive residue row to a primitive integral row and completing it to a determinant-one matrix proves surjectivity. The matrices
\[
 s=\begin{pmatrix}0&-1\\1&0\end{pmatrix},\quad
 r=\begin{pmatrix}0&-1\\1&1\end{pmatrix},\quad
 t=\begin{pmatrix}1&1\\0&1\end{pmatrix}
\]
act on the right, with $s^2=r^3=1$ and $sr=t$ in the projective action. Their cycles give the passport. Their numbers of fixed darts give $e_2,e_3$; the $t$-cycles give the cusp widths. Riemann--Hurwitz gives
\[
 g_0=1+\mu/12-e_2/4-e_3/3-c/2.
\]
The executable \texttt{d\_lambencies.py} constructs every primitive row, canonicalizes its unit orbit, and stores all permutations, rather than only their cycle partitions.

Write $f_d=\gcd(d,N/d)$. The cusp label $c_{d,a}$ denotes the class of $a'/d$, where $d\mid N$, $a\in(\mathbb Z/f_d)^\times$, and $a'\equiv a\pmod{f_d}$ is chosen coprime to $d$. For $f_d=1$ we write $a=1$. In particular $c_{1,1}=0$ and $c_{N,1}=\infty$. A reduced rational cusp $x/y$ has label
\[
 d=\gcd(y,N),\qquad a=x(y/d)\pmod{f_d}.
\]
For $\left(\begin{smallmatrix}A&B\\NC&D\end{smallmatrix}\right)\in\Gamma_0(N)$, its transformed numerator is $Ax$ modulo $f_d$, its denominator divided by $d$ is $D(y/d)$ modulo $f_d$, and $AD\equiv1$. This proves invariance. Conversely, completing the two primitive pairs to unimodular matrices and solving the lower-left congruence by B\'ezout gives equivalence exactly for equal $d,a$; equivalently this is the usual cusp criterion used in the opened modular-curve source \cite{RouseWebb2015}. Its width is obtained from the lower-left entry of $\gamma t^w\gamma^{-1}$: the least positive $w$ with $N\mid d^2w$ is
\[
 w_d=N/\gcd(d^2,N).
\]
The Galois action in \cite[Theorem~16]{RouseWebb2015} sends the numerator invariant to its product with a cyclotomic unit (or its inverse, depending on the chosen action convention). Since that action on $(\mathbb Z/f_d)^\times$ is free and transitive, the exact residue field is $\mathbb Q(\zeta_{f_d})$. This also supplies a direct derivation of the field, rather than merely a field over which the cusp becomes rational.

For an exact divisor $e$, choose integers $b,d_0$ with $ed_0-(N/e)b=1$. The matrix
\[
 W_e=\begin{pmatrix}e&b\\N&ed_0\end{pmatrix},\qquad\det W_e=e,
\]
acts on rational cusps; scaling it by $1/\sqrt e$ gives the standard Atkin--Lehner coset. Reduce the resulting fraction and apply the preceding cusp formula. The program checks all products $W_eW_f=W_{ef/\gcd(e,f)^2}$ on cusps. The denominator changes by $d\mapsto ed/\gcd(e,d)^2$, while the table retains the numerator action as well. The latter matters at $N=18$: $W_9$ exchanges the two cusps at denominator $3$ and the two at denominator $6$. The exact cusp fields on a quotient are fixed fields of the stabilizers of these finite Galois orbit sets. In every nontrivial $S$-quotient in this table all cusps are rational. This includes the two merged conjugate pairs for $18+9$.

D.15 below proves genus zero directly from the eta divisor and proves $S$-invariance. For the added cosets, the normalized trace is $\sqrt e(a+d)$. For $e\geq4$ an elliptic element must have trace zero, hence order two. In the sole $e=3$ case, $N=6$, the determinant equation $3ad-2bc=1$ makes $a,d$ odd, so again a trace of absolute value less than two must be zero. The identity coset has no order-three elliptic points for any of the nine nontrivial $S$-rows, as the explicit dart permutations verify. Thus only order two occurs. Hyperbolic area divided by $|S|$ gives
\[
 \frac{\mu}{6|S|}=-2+c+e_2/2,
\]
which proves every remaining signature entry.
\end{proof}

\begin{result}{D.15}{Proved}{The exact tripod comparison and the eta-product check}
The normalization of $X_0(N)\times_{X(1)}X(2)$ has one component of degree $\mu$ over $X(2)$ for odd $N$, and three components, each of degree $\mu/3$, for even $N$. They are defined over $\mathbb Q$. The component containing the standard cusp is isomorphic to $X_0(4N)$ for odd $N$ and to $X_0(2N)$ for even $N$; the other components differ by permutation of the three tripod punctures. For $4\mid N$ there is also the direct cover $X_0(N)\to X_0(4)$ of degree $\mu/6$.

For all twenty-three rows the exact modular function $T^X=1/\eta_{\pi^X}$ is $S$-invariant and descends to a $\mathbb Q$-coordinate on $X_0(N)/S$, with a unique simple pole and expansion $q^{-1}-d^X+O(q)$. The attached computation supplies the full integral cusp divisor and coefficients through $q^{30}$. The assertions of genus zero and invariance follow from those exact divisors and modular transformation criteria, without an inference from a truncated series.
\end{result}
\begin{proof}
The subgroup $\Gamma(2)$ is generated projectively by $t^2$ and $u_2=\left(\begin{smallmatrix}1&0\\-2&1\end{smallmatrix}\right)$. Its orbits on the dart set are the components of the normalized base change, and the three cusp permutations are $t^2,u_2,(t^2u_2)^{-1}$. For odd $N$, the Chinese remainder theorem makes the image of $\Gamma(2)$ act transitively on primitive rows modulo $N$. For even $N$, reduction modulo $2$ records three primitive row classes, and the same lifting argument makes each class an orbit. All three choices are defined over $\mathbb Q$, since the full two-torsion labeling defining $X(2)$ is defined over $\mathbb Q$. Conjugation by $D=\operatorname{diag}(2,1)$ gives, by direct matrix multiplication,
\[
 D\Gamma_0(4N)D^{-1}=\Gamma(2)\cap\Gamma_0(N)\quad(N\text{ odd}),
\]
\[
 D\Gamma_0(2N)D^{-1}=\Gamma(2)\cap\Gamma_0(N)\quad(N\text{ even}).
\]
This conjugation is the algebraic two-isogeny map with the induced two-torsion labeling, so defines the asserted comparison over $\mathbb Q$. The other components are obtained from the rational $S_3$-action on $X(2)$. For the direct cover use $\Gamma_0(N)\subset\Gamma_0(4)$ and the generators $t,u_4=\left(\begin{smallmatrix}1&0\\-4&1\end{smallmatrix}\right)$. Restricting their dart permutations to the orbit of $(0,1)$ computes that cover and its three-point passport. The direct passports for $N=4,8,12,16$ are respectively
\[
 (1,1,1),\quad(1^2,2,2),\quad(1\,3,1\,3,1\,3),\quad(1^4,4,4).
\]
All component permutations, not just their lengths, are in the JSON certificate.

Let $r_\delta$ be the eta exponents of $T^X$, so $r_\delta=-k_\delta$ for the displayed Frame shape. The opened Newman criterion and Ligozat formula in \cite[p.~2]{RouseWebb2015} apply: the program checks
\[
 \sum r_\delta=0,\quad\sum\delta r_\delta\equiv
 \sum(N/\delta)r_\delta\equiv0\pmod{24},\quad
 \prod\delta^{r_\delta}\in(\mathbb Q^\times)^2.
\]
Thus $T^X$ is a weight-zero function with trivial character on $\Gamma_0(N)$ and no zeros or poles in the upper half-plane. At $c_{d,a}$ its order is the exact rational sum
\[
 \operatorname{ord}_{c_{d,a}}T^X=
 \frac N{24}\sum_{\delta\mid N}
 \frac{\gcd(d,\delta)^2r_\delta}{\gcd(d,N/d)d\delta}.
\]
Every value is integral. For every row the negative part consists exactly of the $|S|$ cusps in the orbit of infinity, each with order $-1$; the whole divisor is invariant under the computed $S$-permutations. Consequently $(T^X\circ W_e)/T^X$ is a nonzero constant.

Here are exact constant-one certificates. For each D-type $N=2e$, the matrix
\[
 \begin{pmatrix}e&-(e+1)/2\\2e&-e\end{pmatrix}
\]
has determinant $e$ and trace zero. Its fixed point lies in the upper half-plane, where $T^X$ is finite and nonzero, forcing the constant to be one. At $(N,e)=(30,6)$ choose the trace-zero matrix $(12,-5;30,-12)$, and at $(30,15)$ choose $(15,-8;30,-15)$; their determinant checks are $6$ and $15$. These generate the group containing $W_{10}=W_6W_{15}$. For $(N,e)=(12,4)$, the cusp $c_{2,1}$ is fixed and has eta order zero, again forcing the constant to be one. These certificates are also checked by the program. Therefore $T^X$ descends to the quotient and has one simple pole there. A nonconstant meromorphic function of degree one identifies that smooth projective quotient with $\mathbb P^1$.

The $q$-product has leading order $-1$ and leading coefficient one; its next coefficient is $k_1=-d^X$, as checked against every root system of the opened CDH Table~2 \cite{CDH2014}. Rational $q$-expansions on the canonical rational modular curve and rationality of the Atkin--Lehner maps give the coordinate over $\mathbb Q$. This independently verifies the content of CDH Proposition~2.1 for the table used here.

The quotient maps retain a precise limitation on these covers. For every nontrivial $e$, the function $j\circ W_e$ has pole order $e$ at infinity, while $j$ has pole order one. Hence $j$ does not factor through $X_0(N)/\langle W_e\rangle$. In the special direct-cover row $12+4$, the map $X_0(12)\to X_0(4)$ sends infinity to infinity but sends $W_4(\infty)=c_{3,1}$ to the cusp zero. It therefore does not factor through that quotient either. The explicit covers above, together with the stated quotient maps, are the resulting tripod form; no descent of those particular maps is asserted.
\end{proof}

The next tables make each cusp field and every nontrivial $S$-action explicit. A set in the second subscript lists the distinct numerator invariants; it is not an unreduced rational fraction. The fields with $f=1,2$ are $\mathbb Q$. In the tripod table the other even-level components have the displayed passport with the three punctures permuted. The JSON and the full readable certificate \texttt{code/d\_lambencies.md} additionally list every quotient orbit and every individual eta order.

{\small\begin{longtable}{r|p{0.79\textwidth}}
$N$ & Cusps $c_{d,a}$, with $(w;f)$ meaning width $w$ and field $\mathbb Q(\zeta_f)$\\\hline
2 & $c_{1,\{1\}}:(2;1)$;\ $c_{2,\{1\}}:(1;1)$\\
3 & $c_{1,\{1\}}:(3;1)$;\ $c_{3,\{1\}}:(1;1)$\\
4 & $c_{1,\{1\}}:(4;1)$;\ $c_{2,\{1\}}:(1;2)$;\ $c_{4,\{1\}}:(1;1)$\\
5 & $c_{1,\{1\}}:(5;1)$;\ $c_{5,\{1\}}:(1;1)$\\
6 & $c_{1,\{1\}}:(6;1)$;\ $c_{2,\{1\}}:(3;1)$;\ $c_{3,\{1\}}:(2;1)$;\ $c_{6,\{1\}}:(1;1)$\\
7 & $c_{1,\{1\}}:(7;1)$;\ $c_{7,\{1\}}:(1;1)$\\
8 & $c_{1,\{1\}}:(8;1)$;\ $c_{2,\{1\}}:(2;2)$;\ $c_{4,\{1\}}:(1;2)$;\ $c_{8,\{1\}}:(1;1)$\\
9 & $c_{1,\{1\}}:(9;1)$;\ $c_{3,\{1,2\}}:(1;3)$;\ $c_{9,\{1\}}:(1;1)$\\
10 & $c_{1,\{1\}}:(10;1)$;\ $c_{2,\{1\}}:(5;1)$;\ $c_{5,\{1\}}:(2;1)$;\ $c_{10,\{1\}}:(1;1)$\\
12 & $c_{1,\{1\}}:(12;1)$;\ $c_{2,\{1\}}:(3;2)$;\ $c_{3,\{1\}}:(4;1)$;\ $c_{4,\{1\}}:(3;1)$;\ $c_{6,\{1\}}:(1;2)$;\ $c_{12,\{1\}}:(1;1)$\\
13 & $c_{1,\{1\}}:(13;1)$;\ $c_{13,\{1\}}:(1;1)$\\
14 & $c_{1,\{1\}}:(14;1)$;\ $c_{2,\{1\}}:(7;1)$;\ $c_{7,\{1\}}:(2;1)$;\ $c_{14,\{1\}}:(1;1)$\\
16 & $c_{1,\{1\}}:(16;1)$;\ $c_{2,\{1\}}:(4;2)$;\ $c_{4,\{1,3\}}:(1;4)$;\ $c_{8,\{1\}}:(1;2)$;\ $c_{16,\{1\}}:(1;1)$\\
18 & $c_{1,\{1\}}:(18;1)$;\ $c_{2,\{1\}}:(9;1)$;\ $c_{3,\{1,2\}}:(2;3)$;\ $c_{6,\{1,2\}}:(1;3)$;\ $c_{9,\{1\}}:(2;1)$;\ $c_{18,\{1\}}:(1;1)$\\
22 & $c_{1,\{1\}}:(22;1)$;\ $c_{2,\{1\}}:(11;1)$;\ $c_{11,\{1\}}:(2;1)$;\ $c_{22,\{1\}}:(1;1)$\\
25 & $c_{1,\{1\}}:(25;1)$;\ $c_{5,\{1,2,3,4\}}:(1;5)$;\ $c_{25,\{1\}}:(1;1)$\\
30 & $c_{1,\{1\}}:(30;1)$;\ $c_{2,\{1\}}:(15;1)$;\ $c_{3,\{1\}}:(10;1)$;\ $c_{5,\{1\}}:(6;1)$;\ $c_{6,\{1\}}:(5;1)$;\ $c_{10,\{1\}}:(3;1)$;\ $c_{15,\{1\}}:(2;1)$;\ $c_{30,\{1\}}:(1;1)$\\
46 & $c_{1,\{1\}}:(46;1)$;\ $c_{2,\{1\}}:(23;1)$;\ $c_{23,\{1\}}:(2;1)$;\ $c_{46,\{1\}}:(1;1)$\\
\end{longtable}
\begin{longtable}{r|p{0.78\textwidth}}
$\ell$ & Atkin--Lehner permutations; omitted cusp labels are fixed\\\hline
$6+3$ & $W_{3}=(c_{1,1}\ c_{3,1})(c_{2,1}\ c_{6,1})$\\
$10+5$ & $W_{5}=(c_{1,1}\ c_{5,1})(c_{2,1}\ c_{10,1})$\\
$14+7$ & $W_{7}=(c_{1,1}\ c_{7,1})(c_{2,1}\ c_{14,1})$\\
$18+9$ & $W_{9}=(c_{1,1}\ c_{9,1})(c_{2,1}\ c_{18,1})(c_{3,1}\ c_{3,2})(c_{6,1}\ c_{6,2})$\\
$22+11$ & $W_{11}=(c_{1,1}\ c_{11,1})(c_{2,1}\ c_{22,1})$\\
$30+15$ & $W_{15}=(c_{1,1}\ c_{15,1})(c_{2,1}\ c_{30,1})(c_{3,1}\ c_{5,1})(c_{6,1}\ c_{10,1})$\\
$46+23$ & $W_{23}=(c_{1,1}\ c_{23,1})(c_{2,1}\ c_{46,1})$\\
$12+4$ & $W_{4}=(c_{1,1}\ c_{4,1})(c_{3,1}\ c_{12,1})$\\
$30+6,10,15$ & $W_{6}=(c_{1,1}\ c_{6,1})(c_{2,1}\ c_{3,1})(c_{5,1}\ c_{30,1})(c_{10,1}\ c_{15,1})$\newline $W_{10}=(c_{1,1}\ c_{10,1})(c_{2,1}\ c_{5,1})(c_{3,1}\ c_{30,1})(c_{6,1}\ c_{15,1})$\newline $W_{15}=(c_{1,1}\ c_{15,1})(c_{2,1}\ c_{30,1})(c_{3,1}\ c_{5,1})(c_{6,1}\ c_{10,1})$\\
\end{longtable}
\begin{longtable}{r|r|r|l|l}
$N$ & components & degree & passport of one component over $X(2)$ & direct degree to $X_0(4)$\\\hline
2 & 3 & 1 & $(1,1,1)$ & --\\
3 & 1 & 4 & $(1\,3,1\,3,1\,3)$ & --\\
4 & 3 & 2 & $(1^{2},2,2)$ & 1\\
5 & 1 & 6 & $(1\,5,1\,5,1\,5)$ & --\\
6 & 3 & 4 & $(1\,3,1\,3,1\,3)$ & --\\
7 & 1 & 8 & $(1\,7,1\,7,1\,7)$ & --\\
8 & 3 & 4 & $(1^{4},4,4)$ & 2\\
9 & 1 & 12 & $(1^{3}\,9,1^{3}\,9,1^{3}\,9)$ & --\\
10 & 3 & 6 & $(1\,5,1\,5,1\,5)$ & --\\
12 & 3 & 8 & $(1^{2}\,3^{2},2\,6,2\,6)$ & 4\\
13 & 1 & 14 & $(1\,13,1\,13,1\,13)$ & --\\
14 & 3 & 8 & $(1\,7,1\,7,1\,7)$ & --\\
16 & 3 & 8 & $(1^{4}\,2^{2},8,8)$ & 4\\
18 & 3 & 12 & $(1^{3}\,9,1^{3}\,9,1^{3}\,9)$ & --\\
22 & 3 & 12 & $(1\,11,1\,11,1\,11)$ & --\\
25 & 1 & 30 & $(1^{5}\,25,1^{5}\,25,1^{5}\,25)$ & --\\
30 & 3 & 24 & $(1\,3\,5\,15,1\,3\,5\,15,1\,3\,5\,15)$ & --\\
46 & 3 & 24 & $(1\,23,1\,23,1\,23)$ & --\\
\end{longtable}
\begin{longtable}{r|l|r|r}
$\ell$ & $\pi^X$ (positive numerator, negative denominator) & $[q^0]T^X$ & $[q]T^X$\\\hline
$2$ & $\frac{2^{24}}{1^{24}}$ & -24 & 276\\
$3$ & $\frac{3^{12}}{1^{12}}$ & -12 & 54\\
$4$ & $\frac{4^{8}}{1^{8}}$ & -8 & 20\\
$5$ & $\frac{5^{6}}{1^{6}}$ & -6 & 9\\
$6$ & $\frac{2\,6^{5}}{1^{5}\,3}$ & -5 & 6\\
$7$ & $\frac{7^{4}}{1^{4}}$ & -4 & 2\\
$8$ & $\frac{2^{2}\,8^{4}}{1^{4}\,4^{2}}$ & -4 & 4\\
$9$ & $\frac{9^{3}}{1^{3}}$ & -3 & 0\\
$10$ & $\frac{2\,10^{3}}{1^{3}\,5}$ & -3 & 1\\
$12$ & $\frac{2^{2}\,3\,12^{3}}{1^{3}\,4\,6^{2}}$ & -3 & 2\\
$13$ & $\frac{13^{2}}{1^{2}}$ & -2 & -1\\
$16$ & $\frac{2\,16^{2}}{1^{2}\,8}$ & -2 & 0\\
$18$ & $\frac{2\,3\,18^{2}}{1^{2}\,6\,9}$ & -2 & 0\\
$25$ & $\frac{25}{1}$ & -1 & -1\\
$6+3$ & $\frac{2^{6}\,6^{6}}{1^{6}\,3^{6}}$ & -6 & 15\\
$10+5$ & $\frac{2^{4}\,10^{4}}{1^{4}\,5^{4}}$ & -4 & 6\\
$14+7$ & $\frac{2^{3}\,14^{3}}{1^{3}\,7^{3}}$ & -3 & 3\\
$18+9$ & $\frac{2^{3}\,3^{2}\,18^{3}}{1^{3}\,6^{2}\,9^{3}}$ & -3 & 3\\
$22+11$ & $\frac{2^{2}\,22^{2}}{1^{2}\,11^{2}}$ & -2 & 1\\
$30+15$ & $\frac{2^{2}\,3\,5\,30^{2}}{1^{2}\,6\,10\,15^{2}}$ & -2 & 1\\
$46+23$ & $\frac{2\,46}{1\,23}$ & -1 & 0\\
$12+4$ & $\frac{2^{4}\,3^{4}\,12^{4}}{1^{4}\,4^{4}\,6^{4}}$ & -4 & 6\\
$30+6,10,15$ & $\frac{2^{3}\,3^{3}\,5^{3}\,30^{3}}{1^{3}\,6^{3}\,10^{3}\,15^{3}}$ & -3 & 3\\
\end{longtable}
}
